
Adversarial corruption is evidence unreliability under an adaptive attacker. RAG systems are vulnerable because they externalise memory into corpora, indices, prompts, and tools. An attacker can target any of these surfaces: the corpus before indexing, the retriever during matching, the prompt template, the retrieved context, or the downstream agent that executes actions.

\begin{table*}[t]
\centering
\small
\begin{tabularx}{\linewidth}{p{0.17\linewidth}p{0.19\linewidth}p{0.25\linewidth}X}
\toprule
\textbf{Attack/defense surface} & \textbf{Representative work} & \textbf{Threat or defense} & \textbf{Implication} \\
\midrule
Retriever poisoning & \citetalias{poisoningRetrievalCorpora2023} & Inject adversarial passages to manipulate dense retrievers & Retrieval security can fail before generation begins. \\
RAG knowledge corruption & \citetalias{poisonedRAG2024} & Craft malicious passages that steer generated answers & End-to-end RAG needs corpus integrity and generation safeguards. \\
Backdoor attacks & \citetalias{trojanRAG2024}; \citetalias{phantom2024} & Triggered retrieval/generation behavior via poisoned contexts & RAG can act as a backdoor carrier without modifying model weights. \\
Prompt template attacks & \citetalias{aip2025} & Weaponize shared instructional prompts to steer retrieval and generation & Prompt templates are part of the attack surface, not just implementation details. \\
Indirect prompt injection & \citetalias{rrag005} & Optimize trigger fragments so malicious instructions are retrieved in the wild & Real attacks must overcome the retrieval barrier, not merely place malicious text in a corpus. \\
Certified defense & \citetalias{certifiedRAG2024} & Isolate-and-aggregate answers for robustness certificates & Certification gives stronger guarantees but typically assumes specific attack budgets. \\
Reranking defense & \citetalias{rrag004} & Use probe-gradient instability to demote poisoned passages & Defense can be inserted retriever-side without retraining the LLM. \\
Attention defense & \citetalias{rrag015} & Detect poisoning via normalised attention variance & Internal attention signals may expose poisoned evidence, but adaptivity remains a concern. \\
\bottomrule
\end{tabularx}
\caption{Representative adversarial corruption attacks and defenses.}
\label{tab:adversarial-methods}
\end{table*}

\begin{table*}[t]
\centering
\scriptsize
\renewcommand{\arraystretch}{0.88}
\begin{tabularx}{\linewidth}{>{\raggedright\arraybackslash}p{0.18\linewidth}c c c c >{\raggedright\arraybackslash}p{0.15\linewidth}X}
\toprule
\textbf{Threat class} & \textbf{Corpus} & \textbf{Query} & \textbf{Retr.} & \textbf{Gen.} & \textbf{Objective} & \textbf{Boundary} \\
\midrule
Dense-retriever poisoning \citep{poisoningRetrievalCorpora2023} & write & train set & white/gray & no & retrieve bad passages & Not always end-to-end. \\
RAG knowledge corruption \citep{poisonedRAG2024,badRAG2024} & write & target topic & black/gray & black & answer steering & Stealth/dynamics vary. \\
Backdoor-style RAG \citep{trojanRAG2024,phantom2024} & write & trigger & gray & black & latent behavior & Normal queries may pass. \\
Indirect prompt injection \citep{aip2025,rrag005} & write/share & natural & black & black & instruction hijack & Must overcome retrieval. \\
Opinion/fairness manipulation \citep{rrag104,rrag162} & write & group/topic & black/gray & black & bias steering & QA metrics miss social harm. \\
Defense setting \citep{certifiedRAG2024,rrag004,rrag015,rrag128} & bounded & varies & assumed & black/gray & detect/demote/trace & Depends on budget/visibility. \\
\bottomrule
\end{tabularx}
\caption{Attacker-capability matrix for adversarial RAG. The same ``attack success rate'' can mean different things depending on corpus write access, query control, retriever visibility, generator access, and objective.}
\label{tab:attacker-capability}
\end{table*}

\begin{table*}[t]
\centering
\scriptsize
\renewcommand{\arraystretch}{0.88}
\begin{tabularx}{\linewidth}{>{\raggedright\arraybackslash}p{0.17\linewidth}>{\raggedright\arraybackslash}p{0.21\linewidth}>{\raggedright\arraybackslash}p{0.22\linewidth}>{\raggedright\arraybackslash}p{0.16\linewidth}X}
\toprule
\textbf{Comparison} & \textbf{Why it looks comparable} & \textbf{Hidden mismatch} & \textbf{Metric pitfall} & \textbf{Normalised reading} \\
\midrule
PoisonedRAG vs. dense-retriever poisoning & Both insert malicious evidence into a retrieval corpus. & One evaluates end-to-end answer steering; the other may evaluate retriever ranking before generation. & Retrieval success and answer attack success are not the same target. & Compare only after recording whether poisoned evidence is retrieved and whether the generator changes the final answer. \\
Certified RAG defense vs. prompt-injection defense & Both are RAG defenses against corrupted retrieved content. & Certification assumes a bounded number of corrupted passages; prompt injection exploits instruction following rather than only false facts. & A certificate under bounded corruption may not imply instruction-hijack resistance. & Treat certification, instruction isolation, and document sanitisation as complementary defenses. \\
Traceback vs. prevention & Both reduce harm after poisoning. & Traceback identifies influential poisoned documents after failure; prevention attempts to block harmful output before response. & Lower post-hoc traceback error does not equal lower user-facing attack success. & Report prevention, detection, demotion, traceback, and recovery as separate outcome variables. \\
Backdoor trigger vs. natural query attack & Both can produce targeted malicious behavior. & Trigger attacks may rely on rare phrases; natural-query attacks must survive ordinary retrieval and ranking. & Attack success under attacker-crafted triggers overstates risk under natural user queries. & Require query-control and trigger-frequency fields before comparing attack strength. \\
\bottomrule
\end{tabularx}
\caption{Concrete examples of threat-model normalisation. These cases show how reported attack or defense success can reverse when attacker access, retrieval realism, and evaluation target differ.}
\label{tab:threat-model-examples}
\end{table*}



\paragraph{Corpus and Retriever Poisoning}

Corpus poisoning attacks insert malicious passages into the retrieval database so that they are returned for target queries. Early work on poisoning retrieval corpora demonstrates that small numbers of adversarial passages can manipulate dense retrieval \citep{poisoningRetrievalCorpora2023}. PoisonedRAG extends this threat to end-to-end generation: retrieved malicious passages can cause the LLM to output attacker-chosen answers \citep{poisonedRAG2024}. BadRAG, UniC-RAG, and related methods explore different trigger structures, optimisation objectives, and attack settings \citep{badRAG2024,rrag049,rrag057}. The important lesson is that retrieval augmentation can amplify small corpus changes into large generation effects.

Backdoor-style attacks make the threat more persistent. TrojanRAG and Phantom show that triggers can activate malicious behavior through retrieved contexts without changing the generator weights \citep{trojanRAG2024,phantom2024}. Other attacks optimise token-level substitutions, influential tokens, or attention-steering patterns \citep{rrag050,rrag052,rrag120}. These works expose a system risk: a RAG system may appear safe under normal queries but fail under targeted concepts, phrases, or document manipulations.

\paragraph{Indirect Prompt Injection and Prompt-Level Attacks}

Indirect prompt injection differs from factual poisoning. Instead of inserting false evidence, the attacker inserts instructions that are retrieved and interpreted by the LLM as commands. In practice, an indirect prompt injection must overcome the retrieval barrier: the malicious document must actually be retrieved under natural queries. \citet{rrag005} explicitly study this barrier and propose trigger fragments that improve retrieval of attack fragments. AIP broadens the surface further by showing that shared instructional prompts can be weaponised to steer RAG retrieval and generation \citep{aip2025}. These attacks imply that prompt templates, document ingestion rules, and retrieved text should be treated as untrusted inputs.

\paragraph{Defenses and Their Assumptions}

Defenses against adversarial corruption operate at multiple layers. Certified approaches such as RobustRAG isolate and aggregate outputs to provide guarantees under bounded retrieval corruption \citep{certifiedRAG2024}. Reranking defenses such as ProGRank use instability or probe-gradient signals to demote suspicious passages \citep{rrag004}. Attention-based filters detect abnormal attention patterns caused by poisoned evidence \citep{rrag015}. Clustering and conflict-resolution defenses such as TrustRAG and SeCon-RAG filter semantically inconsistent or suspicious evidence before generation \citep{rrag013,rrag017}. Traceback methods attempt to identify poisoned texts after an attack has occurred \citep{rrag128}.

The main limitation is threat-model fragmentation. A defense designed for corpus poisoning may not stop prompt-template attacks; a filter trained on obvious malicious instructions may not detect fact-level poisoning; a certified aggregation method may assume a corruption budget that is hard to estimate in open corpora. Robust RAG security therefore needs compositional evaluation: combine poisoned facts, retrieval triggers, prompt-bearing documents, and benign noise in the same pipeline.

\paragraph{Threat Model Dimensions}

Security papers differ in attacker knowledge, access, and objective. Some assume the attacker can insert a few documents into a public corpus but cannot see user queries. Others assume black-box query access, white-box retriever gradients, control over prompt templates, or access to agent memory. Objectives also differ: targeted answer steering, denial of service, refusal induction, opinion manipulation, privacy extraction, fairness manipulation, or guardrail bypass. Comparing attacks without specifying these dimensions is misleading. A poisoning method that works with five injected documents under black-box access has stronger practical relevance than a method requiring full index control; a defense evaluated only against keyword triggers may fail against semantic triggers.

Table~\ref{tab:attacker-capability} makes this comparison explicit. We separate corpus write access, query control, retriever visibility, generator access, and objective because they determine whether two attack results are comparable. This matrix also clarifies why benign retrieval noise and adversarial corruption should not be collapsed into a single robustness score: the former tests system tolerance to accidental evidence errors, while the latter tests adaptive pressure against the retrieval-generation interface.

\paragraph{Attack Success Is Not Enough}

Attack success rate is necessary but insufficient. A strong RAG attack must also be stealthy, durable, and retrieval-realistic. Stealth means poisoned passages should not be trivially detected by lexical filters or human inspection. Durability means attacks should persist under chunking, reranking, context filtering, and prompt changes. Retrieval realism means the malicious content must be retrieved under natural user queries, not only under attacker-crafted queries. Work on indirect prompt injection in the wild explicitly highlights this retrieval barrier \citep{rrag005}. Recent stealth and attention-aware work similarly argues that defenses should evaluate whether poisoned content can be distinguished from benign evidence \citep{rrag015}.

\paragraph{Defense-in-Depth}

The surveyed defenses suggest a defense-in-depth architecture. At ingestion time, systems should validate sources, scan documents for prompt-like instructions, preserve provenance, and detect anomalous clusters. At retrieval time, systems can rerank with robustness signals, compare retrieved documents for contradictions, and demote passages with suspicious trigger behavior. At generation time, systems should isolate evidence from instructions, verify answer support, and refuse to execute instructions found in retrieved documents. During operation, systems need logging and traceback to identify which corpus items influenced harmful outputs. No layer is sufficient alone, but layered controls reduce the chance that one bypass compromises the entire pipeline.
